
\documentclass[aspectratio=169]{beamer}
\usetheme{Madrid}

\usepackage{amsmath,amssymb,amsfonts,bm}
\usepackage{physics}

\title{Quantum Fisher Information and Quantum Cramér–Rao Bound}
\author{Morten Hjorth-Jensen}
\date{March 7, 2026}

\begin{document}

\frame{\titlepage}

\begin{frame}{Structure}
\begin{itemize}
\item Classical Fisher information and Cramér–Rao bound
\item Quantum statistical models
\item Symmetric logarithmic derivative (SLD)
\item Derivation of Quantum Fisher Information (QFI)
\item Quantum Cramér–Rao bound (QCRB)
\item QFI for pure and mixed states
\item Bures metric and fidelity geometry
\item Ramsey interferometry
\item Spin sensors and magnetic field estimation
\item GHZ and squeezed states
\item Multi-parameter estimation
\item Noise and decoherence effects
\item Bayesian estimation
\item Numerical examples
\end{itemize}
\end{frame}

\section{Classical Estimation Theory}

\begin{frame}{Parameter estimation}
Suppose a parameter $\theta$ determines a probability distribution
\begin{equation}
p(x|\theta).
\end{equation}

An estimator $\hat{\theta}$ is constructed from measurement outcomes.

The variance of the estimator is

\begin{equation}
\mathrm{Var}(\hat{\theta})
=
\langle(\hat{\theta}-\theta)^2\rangle .
\end{equation}
\end{frame}

\begin{frame}{Classical Fisher information}
The Fisher information is defined as

\begin{equation}
F(\theta)
=
\sum_x
p(x|\theta)
\left(
\frac{\partial}{\partial \theta}
\ln p(x|\theta)
\right)^2 .
\end{equation}

Equivalent form

\begin{equation}
F(\theta)
=
\sum_x
\frac{1}{p(x|\theta)}
\left(
\frac{\partial p(x|\theta)}{\partial \theta}
\right)^2 .
\end{equation}
\end{frame}

\begin{frame}{Cramér–Rao bound}
For unbiased estimators

\begin{equation}
\mathrm{Var}(\hat{\theta})
\ge
\frac{1}{M F(\theta)}
\end{equation}

where

\begin{itemize}
\item $M$ is number of measurements
\item $F(\theta)$ is Fisher information.
\end{itemize}
\end{frame}

\section{Quantum Estimation}

\begin{frame}{Quantum statistical model}

In quantum mechanics

\begin{equation}
\rho(\theta)
\end{equation}

depends on the unknown parameter.

Measurement outcomes follow

\begin{equation}
p(x|\theta)
=
\mathrm{Tr}[\rho(\theta) E_x].
\end{equation}

where $E_x$ are POVM elements.
\end{frame}

\begin{frame}{Quantum Fisher Information}

The maximal achievable Fisher information over all measurements is

\begin{equation}
F_Q(\theta)
=
\max_{E_x} F(\theta).
\end{equation}

This quantity is the \textbf{Quantum Fisher Information}.
\end{frame}

\section{Symmetric Logarithmic Derivative}

\begin{frame}{SLD definition}

The symmetric logarithmic derivative $L_\theta$ is defined via

\begin{equation}
\frac{\partial \rho}{\partial \theta}
=
\frac12 (\rho L_\theta + L_\theta \rho).
\end{equation}
\end{frame}

\begin{frame}{QFI formula}

The QFI can be written

\begin{equation}
F_Q(\theta)
=
\mathrm{Tr}(\rho L_\theta^2).
\end{equation}
\end{frame}

\section{Pure States}

\begin{frame}{Pure-state QFI}

For

\begin{equation}
\rho = |\psi\rangle \langle \psi|
\end{equation}

the QFI becomes

\begin{equation}
F_Q =
4\left(
\langle \partial_\theta \psi | \partial_\theta \psi \rangle
-
|\langle \psi | \partial_\theta \psi \rangle|^2
\right).
\end{equation}
\end{frame}

\begin{frame}{Unitary encoding}

If

\begin{equation}
|\psi(\theta)\rangle
=
e^{-i\theta G}|\psi_0\rangle
\end{equation}

then

\begin{equation}
F_Q = 4 \mathrm{Var}(G).
\end{equation}
\end{frame}

\section{Spin Sensors}

\begin{frame}{Magnetic field sensing}

Hamiltonian

\begin{equation}
H = \gamma B S_z
\end{equation}

State evolution

\begin{equation}
|\psi(B)\rangle = e^{-i \gamma B t S_z} |\psi_0\rangle
\end{equation}

The QFI becomes

\begin{equation}
F_Q = 4 t^2 \mathrm{Var}(S_z).
\end{equation}
\end{frame}

\section{Ramsey Interferometry}

\begin{frame}{Ramsey protocol}

Steps

\begin{enumerate}
\item Prepare superposition
\item Free evolution
\item Interference measurement
\end{enumerate}

Probability

\begin{equation}
P(\uparrow) = \frac12 (1+\cos(\theta)).
\end{equation}
\end{frame}

\begin{frame}{QFI of Ramsey experiment}

For $N$ independent spins

\begin{equation}
F_Q = N t^2.
\end{equation}

Precision

\begin{equation}
\Delta \theta = \frac{1}{\sqrt{N} t}.
\end{equation}

Standard Quantum Limit.
\end{frame}

\section{Entangled Sensors}

\begin{frame}{GHZ states}

\begin{equation}
|\mathrm{GHZ}\rangle
=
\frac{1}{\sqrt2}(|000\rangle + |111\rangle)
\end{equation}

For $N$ particles

\begin{equation}
F_Q = N^2 t^2.
\end{equation}
\end{frame}

\begin{frame}{Heisenberg limit}

Precision scaling

\begin{equation}
\Delta \theta = \frac{1}{N}.
\end{equation}

Better than standard quantum limit.
\end{frame}

\section{Geometry}

\begin{frame}{Bures metric}

Quantum Fisher information defines the metric

\begin{equation}
ds^2 = \frac14 F_Q d\theta^2.
\end{equation}

This is the Bures metric on quantum state space.
\end{frame}

\section{Noise}

\begin{frame}{Decoherence}

Dephasing reduces QFI

\begin{equation}
F_Q(t) = N t^2 e^{-\gamma t}.
\end{equation}

Optimal sensing time exists.
\end{frame}

\section{Bayesian Estimation}

\begin{frame}{Bayesian approach}

Posterior distribution

\begin{equation}
p(\theta|x)
\propto
p(x|\theta)p(\theta).
\end{equation}

Allows sequential updating.
\end{frame}

\begin{frame}{Summary}

Key results:

\begin{itemize}
\item Fisher information quantifies sensitivity
\item QFI gives optimal achievable precision
\item Quantum Cramér–Rao bound

\begin{equation}
\mathrm{Var}(\hat{\theta})
\ge
\frac{1}{M F_Q}.
\end{equation}

\item Entanglement enhances precision
\end{itemize}

\end{frame}

\end{document}
